\chapter{Neck Stiffness}

\section*{Definition}
Neck stiffness is difficulty or pain when moving the neck, often accompanied by a feeling of tightness. It ranges from a benign muscle strain to a warning sign of serious illness: when stiffness comes on quickly with fever, severe headache, or confusion, it can signal meningitis or bleeding around the brain.

\section*{When to See a Doctor}
Get emergency care if neck stiffness is accompanied by any of these:
\begin{itemize}
  \item Fever, especially with a severe or worsening headache
  \item Confusion, drowsiness, or personality change
  \item Sensitivity to light, nausea and vomiting, or a rash that does not fade when pressed
  \item A sudden, explosive "worst headache of your life"
  \item Seizures, weakness, or numbness, or stiffness after a head or neck injury
\end{itemize}
Without these warning signs, see your doctor if stiffness lasts more than a week or keeps recurring.

\section*{How It Develops}
In a benign strain, the neck muscles go into spasm -- usually after poor sleep posture, prolonged screen work, or a minor injury -- and the tightness restricts movement until the spasm relaxes. In meningitis, however, the meninges, the three thin membranes that wrap the brain and spinal cord, become inflamed by infection. Because these membranes are anchored to the skull and spine, bending the neck forward stretches them; the resulting pain and involuntary resistance is called nuchal rigidity, and it is the body's way of bracing the head. The same protective tightening can occur when blood from a ruptured vessel irritates the meninges, as in subarachnoid hemorrhage. The distinction matters: a strained muscle feels tight but does not cause fever or confusion, whereas meningeal irritation often does.

\section*{Causes}
\begin{itemize}
  \item Benign: Muscle strain from poor posture or awkward sleeping positions, cervical spondylosis (wear-and-tear arthritis of the neck), and whiplash after a minor collision
  \item Serious: Bacterial meningitis (a medical emergency), viral meningitis, encephalitis, and cerebral abscess
  \item Vascular: Subarachnoid hemorrhage from a ruptured brain aneurysm, described as the "worst headache of your life"
  \item Structural: Cervical disc disease and, after trauma, injury to the cervical spine
  \item Other: Epidural or retropharyngeal abscess, and rarely tumors or inflammatory disease
\end{itemize}

\section*{Related Symptoms}
\begin{itemize}
  \item Headache, often severe and generalized
  \item Sensitivity to bright light (photophobia)
  \item Fever, chills, nausea, and vomiting
  \item Pain radiating into the shoulders or between the shoulder blades
  \item Confusion, drowsiness, or difficulty concentrating
  \item A skin rash in some forms of meningitis
\end{itemize}

\section*{Medical Evaluation}
\begin{itemize}
  \item Your doctor asks when the stiffness began, whether it is painful, and whether fever, headache, or other symptoms are present.
  \item The examination checks neck range of motion and looks for nuchal rigidity by gently bending the chin toward the chest; Kernig's and Brudzinski's signs are checked when meningitis is suspected.
  \item A full neurologic examination looks for weakness, numbness, abnormal reflexes, or changes in alertness.
  \item If meningitis is suspected, blood cultures are drawn and blood tests check for inflammation; a lumbar puncture (spinal tap) analyzes the fluid around the brain and spinal cord for infection or blood.
  \item CT or MRI of the brain and spine detect bleeding, abscess, disc disease, or tumors, and CT is done before the spinal tap if there is any concern about pressure in the skull.
\end{itemize}

\section*{Treatment Options}
\begin{itemize}
  \item Bacterial meningitis is treated in the hospital with intravenous antibiotics given immediately, often alongside dexamethasone to reduce inflammation and complications.
  \item Viral meningitis usually resolves with rest, fluids, and pain relief; antiviral drugs are used when a specific virus is identified.
  \item Subarachnoid hemorrhage is treated by a neurosurgeon or interventional neuroradiologist, often by clipping or coiling the aneurysm, plus blood-pressure control.
  \item Benign muscle strain responds to NSAIDs, heat, gentle stretching, and physical therapy.
  \item Cervical arthritis is managed with pain relievers, posture correction, and physiotherapy; severe disc problems may need surgery.
  \item Close monitoring is essential because conditions that cause meningeal irritation can worsen quickly.
\end{itemize}


\section*{Self-Care and Home Management}
\begin{itemize}
  \item For a benign muscle strain, apply a warm compress or heating pad for 15--20 minutes and gently stretch the neck in all directions.
  \item Take acetaminophen, ibuprofen, or naproxen for pain and spasm, following label instructions, unless you cannot take NSAIDs.
  \item Rest in a position that supports your neck, use a supportive pillow, and avoid carrying heavy bags on one shoulder.
  \item Improve posture at your desk: keep the screen at eye level and take breaks from screen work every 30--60 minutes.
  \item Never try to "crack" the neck forcefully, and never ignore stiffness combined with fever, headache, or confusion -- that combination needs emergency evaluation, not home treatment.
\end{itemize}

\section*{References}
\begin{thebibliography}{99}
\bibitem{neck_stiffness:meningitis} Tunkel AR, Hartman BJ, Kaplan SL, et al. Practice guidelines for the management of bacterial meningitis. \textit{Clin Infect Dis}. 2004;39:1267--1284.
\bibitem{neck_stiffness:nejm-meningitis} van de Beek D, de Gans J, Tunkel AR, et al. Community-acquired bacterial meningitis in adults. \textit{N Engl J Med}. 2006;354:44--53.
\bibitem{neck_stiffness:idsa-meningitis} Hasbun R, Rosenthal N, Balada-Llasat JM, et al. Clinical practice guideline for the evaluation and management of suspected bacterial meningitis in adults. \textit{Clin Infect Dis}. 2021;73:e59--e77.
\bibitem{neck_stiffness:escmid} van de Beek D, Cabellos C, Dzupova O, et al. ESCMID guideline: diagnosis and treatment of acute bacterial meningitis. \textit{Clin Microbiol Infect}. 2016;22(suppl 3):S37--S62.
\end{thebibliography}
